\section{Zadanie 4}
Udowodnij, że wskończonym łańcuchu Markova każdu stan powracający jest dodatni.

\subsection*{Dowód}
Niech \(S\) będzie przestrzenią stanów łańcucha Markowa, przy czym \(|S| < \infty\) (przestrzeń jest skończona).
Rozważmy stan powracający \(i \in S\). Niech \(C\) będzie klasą komunikacji, do której należy stan \(i\).

Z własności łańcuchów Markowa wiemy, że:
\begin{enumerate}
	\item Klasa \(C\) jest zamknięta (prawdopodobieństwo wyjścia z niej wynosi 0).
	\item Wszystkie stany w klasie \(C\) są powracające.
	\item Ponieważ cała przestrzeń \(S\) jest skończona, to \(C \subseteq S\) również jest zbiorem skończonym.
\end{enumerate}

Przeprowadzimy dowód nie wprost. \\
Załóżmy przeciwnie, że stan \(i\) (a zatem i wszystkie stany w klasie \(C\)) jest powracający zerowy.

Z definicji stanu powracającego zerowego wynika,
że granica prawdopodobieństw przejścia dąży do zera, czyli dla każdego \(j \in C\):
\[
	\lim_{n \to \infty} P_{ij}^{(n)} = 0
\]
gdzie \(P_{ij}^{(n)}\) oznacza prawdopodobieństwo przejścia ze stanu \(i\) do stanu \(j\) w \(n\) krokach.

Ponieważ klasa \(C\) jest zamknięta, w każdym kroku \(n\) suma prawdopodobieństw przejścia ze stanu \(i\) do wszystkich stanów tej klasy musi wynosić 1:
\[
	\sum_{j \in C} P_{ij}^{(n)} = 1
\]

Obłóżmy to granicą:
\[ \lim_{n \to \infty} \sum_{j \in C} P_{ij}^{(n)} = 1 \]

Ponieważ zbiór \(C\) jest skończony,
możemy zamienić kolejność sumowania i granicy (kluczowe)
\[ \sum_{j \in C} \left( \lim_{n \to \infty} P_{ij}^{(n)} \right) = 1 \]

Korzystając z naszego założenia podstawiamy \(0\) w miejsce granicy:
\[ \sum_{j \in C} 0 = 1 \]

Otrzymaliśmy sprzeczność
Zatem, stan \(i\) musi być powracający dodatni. \ \  \(\square\)
